\begin{proofbox}
\textbf{证明一：最快证明（首选，三步，直接调用切线判定式）。}

设经过 $M(x_0,y_0)$ 的\textbf{非竖直}切线为 $y=kx+m$。
\begin{align*}
&\text{经过 }M:\quad m=y_0-kx_0;\\
&\text{与椭圆相切}:\quad m^2=a^2k^2+b^2;\\
&\text{直接代入}:\quad (y_0-kx_0)^2=a^2k^2+b^2.
\end{align*}
展开后即
\[
\boxed{(a^2-x_0^2)k^2+2x_0y_0k+b^2-y_0^2=0.}
\]
\textbf{速度秘诀：}把“$m=y_0-kx_0$”塞进“$m^2=a^2k^2+b^2$”，不必重新联立椭圆。

\medskip
\textbf{证明二：纯代数通法（联立直线与椭圆，判别式为零）。}

将 $y=kx+m$ 代入 $\frac{x^2}{a^2}+\frac{y^2}{b^2}=1$，两边乘 $a^2b^2$：
\[
(b^2+a^2k^2)x^2+2a^2kmx+a^2(m^2-b^2)=0.
\]
交点横坐标只有一个实根，当且仅当
\begin{align*}
\Delta_x
&=(2a^2km)^2-4(b^2+a^2k^2)a^2(m^2-b^2)\\
&=4a^2b^2(a^2k^2+b^2-m^2)=0.
\end{align*}
因 $a,b>0$，故 $m^2=a^2k^2+b^2$；再代入 $m=y_0-kx_0$，即得 (C009-1)。此证明不预设 C007，属于可独立使用的通法。

\medskip
\textbf{证明三：设切点法（几何与代数结合，也能直接找到切点）。}

设切点为 $T(u,v)$，椭圆在 $T$ 处的切线为
\[
\frac{ux}{a^2}+\frac{vy}{b^2}=1.
\]
若该切线非竖直，则 $v\ne0$。化成斜截式：
\[
y=-\frac{b^2u}{a^2v}x+\frac{b^2}{v}=kx+m.
\]
所以 $u=-\frac{a^2k}{m},\ v=\frac{b^2}{m}$（这里 $m\ne0$，因为切线不经过椭圆中心）。
把切点代入椭圆：
\[
1=\frac{u^2}{a^2}+\frac{v^2}{b^2}
=\frac{a^2k^2+b^2}{m^2}
\Longrightarrow m^2=a^2k^2+b^2.
\]
又因切线经过 $M$，有 $m=y_0-kx_0$，故得到 (C009-1)。
\textbf{顺便得到切点坐标：}
\[
\boxed{T\left(-\frac{a^2k}{m},\frac{b^2}{m}\right),\quad m=y_0-kx_0.}
\]

\medskip
\textbf{证明四：三角参数与柯西不等式（最直观的数形结合）。}

取椭圆动点 $P(a\cos t,b\sin t)$。固定 $k$，令
\[
F(t)=y-kx=b\sin t-ak\cos t.
\]
由柯西不等式或辅助角公式，
\[
|F(t)|\le\sqrt{(-ak)^2+b^2}\sqrt{\cos^2t+\sin^2t}
=\sqrt{a^2k^2+b^2}.
\]
两个端点值均能取到。直线 $y-kx=m$ 要与椭圆恰好相切，必须且只需
\[
m=\pm\sqrt{a^2k^2+b^2}.
\]
故任意 $k\in\mathbb R$ 都对应两条平行切线，且把 $m=y_0-kx_0$ 代入又得到 (C009-1)。

\medskip
\textbf{证明五：化椭圆为圆（几何证明切线存在的条数）。}

作可逆线性变换
\[
U=\frac{x}{a},\qquad V=\frac{y}{b}.
\]
椭圆变成 $U^2+V^2=1$，$M$ 变成 $M'=(x_0/a,y_0/b)$。可逆线性变换将切线一一映成切线，因此
\[
\begin{array}{c|c|c}
OM'^2<1 & OM'^2=1 & OM'^2>1\\ \hline
0\text{ 条} &1\text{ 条}&2\text{ 条}
\end{array}
\]
分别对应椭圆内、椭圆上、椭圆外。注意上述判定计数的是\textbf{切线条数}，并非实数斜率的个数。

\medskip
\textbf{证明六：判别式与韦达定理（最快判断两个有限斜率）。}

记 $S=\frac{x_0^2}{a^2}+\frac{y_0^2}{b^2}$。对 (C009-1)，
\begin{align*}
\Delta_k
&=(2x_0y_0)^2-4(a^2-x_0^2)(b^2-y_0^2)\\
&=4(a^2y_0^2+b^2x_0^2-a^2b^2)\\
&=\boxed{4a^2b^2(S-1)}.
\end{align*}
当 $a^2-x_0^2\ne0$ 时，$\Delta_k$ 的正、零、负分别对应两个、一个、零个\textbf{实数斜率}；两个不同实根时，
\[
k_1+k_2=-\frac{2x_0y_0}{a^2-x_0^2},\qquad
k_1k_2=\frac{b^2-y_0^2}{a^2-x_0^2}.
\]
若 $x_0=\pm a$，须改用竖直切线 $x=x_0$ 单独补齐。$y_0\ne0$ 时另一条斜切线由 $2x_0y_0k+b^2-y_0^2=0$ 给出；$y_0=0$ 时唯一切线是 $x=x_0$。

\medskip
\textbf{补充：椭圆上某点的切线斜率为什么唯一？}

设 $P(x_0,y_0)$ 在椭圆上。隐函数求导可得
\[
\frac{2x}{a^2}+\frac{2y}{b^2}y'=0
\quad\Longrightarrow\quad
\boxed{k=-\frac{b^2x_0}{a^2y_0}}\qquad(y_0\ne0).
\]
它与 (C009-1) 的重根 $k=-\frac{x_0y_0}{a^2-x_0^2}$ 一致。若 $y_0=0$，则 $P=(\pm a,0)$，唯一切线竖直，没有有限斜率。
\end{proofbox}
